\subsection{Una equivalencia notable de categorías.}
\label{subsection:IV.18.1}

\begin{proposition}[18.1.1]
\label{IV.18.1.1}
Sean $S$ un preesquema, $S_0$ un subpreesquema cerrado de $S$, $X_0$ un $S_0$-preesquema liso (resp. étale) sobre $S_0$, y $x_0$ un punto de $X_0$. Existen entonces un entorno abierto $U_0$ de $x_0$ en $X_0$, un $S$-preesquema $U$ liso (resp. étale) sobre $S$ y un $S_0$-isomorfismo
\[
  U\times_S S_0\xrightarrow{\sim}U_0.
\]
\end{proposition}

Observemos que si $X_0$ es étale sobre $S_0$ en el punto $x_0$, es con mayor razón no ramificado sobre $S_0$ en dicho punto; si se ha construido un $S$-preesquema $U$ liso sobre $S$ tal que $U\times_S S_0$ sea isomorfo a $U_0$, como las fibras de los morfismos $U_0\to S_0$ y $U\to S$ que contienen a $x_0$ son entonces isomorfas, se seguirá que $U$ es no ramificado sobre $S$ en el punto $x_0$ (17.4.1, $d$)), y por tanto también en un entorno de $x_0$; sustituyendo $U$ por este entorno, se concluye que $U$ será étale sobre $S$. Basta, pues, demostrar la proposición cuando solo se supone que $X_0$ es liso sobre $S_0$.

Como la cuestión es local sobre $S$ y sobre $X_0$, se puede suponer que $S=\Spec(A)$ y $X_0=\Spec(C_0)$ son afines, de modo que $S_0=\Spec(A_0)$, donde $A_0$ es un anillo cociente de $A$, $C_0=B_0/\mathfrak J_0$, donde $B_0=A_0[T_1,\ldots,T_n]$ y $\mathfrak J_0$ es un ideal de tipo finito de $B_0$; por último, $C_0$ es una $A_0$-álgebra formalmente lisa para las topologías discretas. Sea $\mathfrak p_0$ el ideal $\mathfrak j_{x_0}$ en $C_0$; se tiene $\mathfrak p_0=\mathfrak q_0/\mathfrak J_0$, donde $\mathfrak q_0$ es un ideal primo de $B_0$. El criterio jacobiano (0, 22.6.4), junto con (0, 19.1.12), muestra que existen en $\mathfrak J_0$ una familia de $r$ polinomios $u_i$ $(1\leq i\leq r)$ y $r$ índices $j_h$ $(1\leq h\leq r)$ tales que las imágenes de los $u_i$ en $(\mathfrak J_0)_{\mathfrak q_0}/(\mathfrak J_0^2)_{\mathfrak q_0}$ generan este $(B_0)_{\mathfrak q_0}$-módulo y se tiene
\begin{equation}
  \det\left(\frac{\partial u_i}{\partial T_{j_h}}\right)\notin\mathfrak q_0.
  \tag{18.1.1.1}\label{IV.18.1.1.1}
\end{equation}

Como $(B_0)_{\mathfrak q_0}$ es un anillo local, se sigue del lema de Nakayama que se puede suponer que las imágenes de los $u_i$ en $(\mathfrak J_0)_{\mathfrak q_0}$ generan este $(B_0)_{\mathfrak q_0}$-módulo y, después, sustituyendo si es preciso $X_0$ por un entorno abierto afín de $x_0$, que los $u_i$ generan $\mathfrak J_0$ ($0_{\mathrm I}$, 5.2.2). Pongamos entonces $B=A[T_1,\ldots,T_n]$; $B_0$ es, por tanto, un anillo cociente

\oldpage[IV-4]{110}

de $B$, $\mathfrak q_0$ es la imagen de un ideal primo $\mathfrak q$ de $B$ y $\mathfrak q$ es la imagen inversa de $\mathfrak q_0$. Para todo $i$, sea $v_i\in B$ un elemento del que $u_i$ es la imagen en $B_0$, y sea $\mathfrak J$ el ideal de $B$ generado por los $v_i$, de modo que $\mathfrak J_0$ es la imagen de $\mathfrak J$ en $B_0$. La proposición quedará demostrada tomando como $U$ un entorno abierto del punto de $\Spec(B/\mathfrak J)$ correspondiente al ideal primo $\mathfrak p=\mathfrak q/\mathfrak J$, con tal de que se demuestre que $B_{\mathfrak q}/\mathfrak J_{\mathfrak q}$ es una $A$-álgebra formalmente lisa para las topologías discretas. Ahora bien, esto se sigue del criterio jacobiano, pues las imágenes de los $v_i$ en $\mathfrak J_{\mathfrak q}/\mathfrak J_{\mathfrak q}^2$ generan este $B_{\mathfrak q}$-módulo, y de (18.1.1.1) se sigue que se tiene $\det(\partial v_i/\partial T_{j_h})\notin\mathfrak q$.

\begin{theorem}[18.1.2]
\label{IV.18.1.2}
Sean $S$ un preesquema y $S_0$ un subpreesquema cerrado de $S$ cuyo espacio subyacente es idéntico al de $S$. Entonces el funtor
\[
  X\longmapsto X\times_S S_0
\]
de la categoría de los $S$-preesquemas étales sobre $S$ a la categoría de los $S_0$-preesquemas étales sobre $S_0$ es una equivalencia de categorías.
\end{theorem}

Demostremos primero que este funtor es plenamente fiel. Sean $X$, $Y$ dos $S$-preesquemas étales sobre $S$, y pongamos $X_0=X\times_S S_0$, $Y_0=Y\times_S S_0$. Si $Z=X\times_S Y$, el conjunto $\Hom_S(X,Y)$ está en correspondencia biyectiva canónica con el conjunto de las $X$-secciones $\Gamma(Z/X)$, y del mismo modo $\Hom_{S_0}(X_0,Y_0)$ está en correspondencia biyectiva canónica con $\Gamma(Z_0/X_0)$, donde $Z_0=Z\times_S S_0=X_0\times_{S_0}Y_0$. Ahora bien, $Z$ es étale sobre $X$, $Z_0$ es étale sobre $X_0$, y $X_0$ (resp. $Z_0$) es un subpreesquema cerrado de $X$ (resp. $Z$) con el mismo espacio subyacente. Los abiertos de $Z$ tales que la restricción del morfismo $Z\to X$ sea sobreyectiva y radicial son, por tanto, los mismos que los abiertos de $Z_0$ que tienen las propiedades correspondientes, y nuestra afirmación se sigue así de (17.9.3).

Para terminar la demostración, basta ver que para todo $S_0$-preesquema $X_0$ étale sobre $S_0$, existe un $S$-preesquema $X$ étale sobre $S$ y un $S_0$-isomorfismo $X_0\xrightarrow{\sim}X\times_S S_0$. En virtud de la prop. (18.1.1), existe un recubrimiento abierto $(U_\alpha)$ de $X_0$ y, para todo $\alpha$, un $S$-preesquema $V_\alpha$ que es étale sobre $S$, y finalmente un $S_0$-isomorfismo $\theta_\alpha:U_\alpha\xrightarrow{\sim}V_\alpha\times_S S_0$. Además, en virtud de la primera parte de la demostración, existe un único $S_0$-isomorfismo $\varphi_{\alpha\beta}$ de $\theta_\alpha(U_\alpha\cap U_\beta)$ sobre $\theta_\beta(U_\alpha\cap U_\beta)$, correspondiente al automorfismo identidad de $U_\alpha\cap U_\beta$, y es inmediato, por la misma razón, que estos isomorfismos satisfacen la condición de pegado ($0_{\mathrm I}$, 4.1.7). Existe, por tanto, un $S$-preesquema $X$ tal que los $V_\alpha$ se identifican canónicamente con subpreesquemas inducidos sobre abiertos de $X$, y los $\theta_\alpha$ se identifican con $S_0$-isomorfismos que coinciden en las intersecciones $U_\alpha\cap U_\beta$ y definen, por consiguiente, un $S_0$-isomorfismo $X_0\xrightarrow{\sim}X\times_S S_0$. Es claro que $X$ es étale sobre $S$ (17.3.2), lo que termina la demostración.

\begin{corollary}[18.1.3]
\label{IV.18.1.3}
Sean $S$ un preesquema, $X$ un $S$-preesquema étale sobre $S$, $S'$ un $S$-preesquema, y $S'_0$ un subpreesquema cerrado de $S'$ con el mismo espacio subyacente. Entonces la aplicación canónica $X(S')_S\to X(S'_0)_S$ (I, 3.4.3) es biyectiva.
\end{corollary}

En efecto, pongamos $X'=X\times_S S'$, $X'_0=X\times_S S'_0$, de modo que $X(S')_S=\Gamma(X'/S')$ y $X(S'_0)_S=\Gamma(X'_0/S'_0)$; el corolario se sigue del hecho de que el funtor definido en (18.1.2) (sustituyendo $S$ y $S_0$ por $S'$ y $S'_0$, respectivamente) es plenamente fiel (o directamente de (17.9.3)).
